% Candidate rules of the credential data model that have no Lean yet (§4.1).
% Written by hand; see docs/one-id-survey.md.  The section heading is in credentials.tex.

\begin{definition}[The security model of the photoid (C-26)]
  \label{def:c26}
  \uses{def:element, def:photoid}
  The security model of ISO/IEC TS 23220-4 Annexes A and B must be followed:
  an issuer-signed security object, device-signed responses, and selective
  disclosure of individual data elements. The guide says these ensure a
  verifier receives only what it requested and the holder consented to, with
  authenticity and integrity.

  The sentence states guarantee G-01 as a consequence of the ISO mechanisms.
  The model has to derive it from the checks the guide lists, not assume it.

  External: ISO/IEC TS 23220-4:2026 Annex A, Annex B.

  Source:
  \href{https://standards.atlassian.net/wiki/spaces/OneID/pages/2947416082/Technical+Implementation+Guide}{One
  ID guide} §4.1.1.
\end{definition}

\begin{definition}[Age attestation (C-31)]
  \label{def:c31}
  \uses{def:element, def:photoid, def:c04, def:c103}
  The base identity namespace includes age attestation elements, which
  disclose an age limit without the age or the date of birth when a relying
  party requests them alone.

  An age attestation is a data element like any other, so the rule for
  addressing one element covers it. The guarantee that the date of birth is
  not disclosed with it belongs to G-01.

  External: ISO/IEC TS 23220-4:2026 Annex C, Table C.1.

  Source:
  \href{https://standards.atlassian.net/wiki/spaces/OneID/pages/2947416082/Technical+Implementation+Guide}{One
  ID guide} §4.1.1.
\end{definition}

\begin{definition}[Showing the boarding pass in the wallet (C-39)]
  \label{def:c39}
  \uses{def:boarding-pass, def:c05, def:c102}
  \notready
  The wallet must display 'only all mandatory elements' of the boarding pass
  string in a user-friendly way.

  Needs the list of mandatory elements from Resolution 792. Not modelled.

  Open point: Does the wallet show all mandatory elements and nothing else,
  or at least the mandatory ones?

  External: IATA Resolution 792.

  Source:
  \href{https://standards.atlassian.net/wiki/spaces/OneID/pages/2947416082/Technical+Implementation+Guide}{One
  ID guide} §4.1.2, consideration 2.
\end{definition}

\begin{definition}[How many boarding pass credentials a journey has (C-42)]
  \label{def:c42}
  \uses{def:boarding-pass, def:c03, def:c102}
  \notready
  The guide speaks of the last leg of a journey, so one boarding pass
  credential may cover several flights. It does not say whether a journey has
  one credential or one per leg, or what a wallet holds for several
  passengers.

  A request names a document type. If a wallet holds several boarding passes,
  the rule that picks the one to send is missing.

  Open point: Is there one boarding pass credential per journey, per flight
  or per boarding pass string? When a wallet holds several, which one answers
  a request for the document type?

  Source:
  \href{https://standards.atlassian.net/wiki/spaces/OneID/pages/2947416082/Technical+Implementation+Guide}{One
  ID guide} §4.1.2.
\end{definition}

\begin{definition}[Who issues a boarding pass (C-43)]
  \label{def:c43}
  \uses{def:boarding-pass, def:c01, def:c15}
  \notready
  The guide says a boarding pass is signed when it is issued and does not say
  who issues it or how a relying party trusts that issuer. Its trust wording
  is about issuing authority certificate lists.

  The issuer check of a boarding pass needs a list of boarding pass issuers,
  which may differ from the list of photoid issuers.

  Open point: Does each airline sign its own boarding passes, and how does an
  airport trust the airlines that fly there?

  Source:
  \href{https://standards.atlassian.net/wiki/spaces/OneID/pages/2947416082/Technical+Implementation+Guide}{One
  ID guide} §4.1.2.
\end{definition}

\begin{definition}[Which face image a touchpoint matches (C-97)]
  \label{def:c97}
  \uses{def:photoid, def:photoid-advisories}
  \notready
  A photoid may carry two face images: the portrait, which holds the
  passport's image, and the image captured at enrolment, which the guide says
  may be newer and of higher quality. The guide does not say which one a
  relying party matches a traveler's face against.

  A request names the data elements it wants, so the choice shows in the
  request. No rule says which is right.

  Open point: When a photoid carries both, which image does a relying party
  request and match: the portrait, the image captured at enrolment, or
  either?

  Source:
  \href{https://standards.atlassian.net/wiki/spaces/OneID/pages/2947416082/Technical+Implementation+Guide}{One
  ID guide} §4.1.1.
\end{definition}

\begin{definition}[A start of validity (C-98)]
  \label{def:c98}
  \uses{def:expired}
  \notready
  The guide names the end of a credential's validity and says nothing of a
  start.

  The model has no start of validity, so a credential is valid from any time
  until its validity ends.

  Open point: Does a credential have a start of validity that a relying party
  checks?

  Lookup: ISO/IEC 18013-5:2021 §9.3.1 and the validity information set when
  an mdoc is signed.

  Source:
  \href{https://standards.atlassian.net/wiki/spaces/OneID/pages/2947416082/Technical+Implementation+Guide}{One
  ID guide} §4.1.2.
\end{definition}

\begin{definition}[What a boarding pass string contains (C-102)]
  \label{def:c102}
  \uses{def:boarding-pass}
  \notready
  The guide refers the format of the boarding pass string, and which of its
  elements are mandatory, to IATA Resolution 792. It reproduces none of it.
  IATA's public guide to the resolution was read. It says one string holds
  one flight leg or several, that its items are mandatory, conditional or for
  an airline's own use, and that the passenger name is a mandatory item (BCBP
  guide §2.2.1). It leaves the list of items, their sizes and their values to
  the resolution's Attachments A and C (BCBP guide §2.2.1, §2.2.2).

  The model takes what the resolution decides as a parameter. With the
  resolution read, the passenger name, the legs and the mandatory elements of
  a string become definitions.

  Open point: Which items of a boarding pass string are mandatory, how long is
  the passenger name item and how is a longer name cut, and how are the legs
  of a journey and their dates given?

  Lookup: IATA Resolution 792, Attachments A and C.

  External: IATA Resolution 792 (not read); IATA BCBP Implementation Guide,
  seventh edition, §2.2.1, §2.2.2 (read).

  Source:
  \href{https://standards.atlassian.net/wiki/spaces/OneID/pages/2947416082/Technical+Implementation+Guide}{One
  ID guide} §4.1.2.
\end{definition}

\begin{definition}[What Tables C.1 to C.3 list (C-103)]
  \label{def:c103}
  \uses{def:photoid}
  \notready
  The guide refers the identifier, description, presence and encoding of
  every photoid data element to three tables of ISO/IEC TS 23220-4. It
  reproduces none of them.

  The model takes the tables as a parameter. With the tables read, they
  become one definition, and each guarantee proven for every table holds for
  it.

  Open point: Which data elements do the three tables list, which do they
  mark mandatory, and in which encoding?

  Lookup: ISO/IEC TS 23220-4:2026 Annex C, Tables C.1 to C.3.

  External: ISO/IEC TS 23220-4:2026 Annex C, Tables C.1 to C.3.

  Source:
  \href{https://standards.atlassian.net/wiki/spaces/OneID/pages/2947416082/Technical+Implementation+Guide}{One
  ID guide} §4.1.1.
\end{definition}

\begin{definition}[A boarding pass that covers an adult and an infant (C-104)]
  \label{def:c104}
  \uses{def:boarding-pass, def:traveler-match}
  \notready
  IATA's public guide to the bar code says some airlines issue one boarding
  pass for an adult travelling with an infant (BCBP guide §2.3.5). The One ID
  guide does not say whether a boarding pass credential may cover two
  travelers.

  The model gives a credential one subject. If one boarding pass covers two
  travelers, 'the same traveler' is not defined for it.

  Open point: May a boarding pass credential cover an adult and an infant? If
  it may, which of the two does it refer to when it is matched with a photoid,
  and what does the infant present?

  Blocks: G-03 (a boarding pass with two travelers has no one traveler to
  match).

  External: IATA BCBP Implementation Guide, seventh edition, §2.3.5 (read).

  Source:
  \href{https://standards.atlassian.net/wiki/spaces/OneID/pages/2947416082/Technical+Implementation+Guide}{One
  ID guide} §4.1.2.
\end{definition}

\begin{definition}[A relying party's check of a boarding pass's validity against the journey (C-109)]
  \label{def:c109}
  \uses{def:bp-validity}
  \notready
  The rule that a boarding pass's validity ends on or soon after the
  departure of the last leg is a duty of the issuer (C-38). No rule tells a
  relying party to check that the validity fits the journey. The only time
  check the guide gives a relying party is that a credential is within its
  validity (C-44).

  OID-104 is stated for an issuer. The check a relying party makes does not
  read which departure is meant, so for one credential its answer is the same
  under both readings. In the model the departure is an input of the issuer's
  rule and is not in the credential, so a relying party's check has nothing
  to compare the validity with.

  Case to check: A flight scheduled for 10:00 departs at 13:00, and the
  airline signs the boarding pass again so that its validity ends at 14:00.
  With ``soon after'' read as an hour, the validity breaks no requirement if
  the actual departure is meant, and ends long after the departure if the
  scheduled one is. At 12:30 the boarding pass is not expired under either
  reading, so a relying party that checks expiry accepts it in both.

  Open point: Must a relying party refuse a boarding pass whose validity does
  not fit the journey, or is expiry the only check?

  Source:
  \href{https://standards.atlassian.net/wiki/spaces/OneID/pages/2947416082/Technical+Implementation+Guide}{One
  ID guide} §4.1.2, consideration 1.
\end{definition}
